\bookchapter{第二十五章：二手手機處理與數字遺產}

\chapterintro{嘿，朋友哋，經過前面二十幾章嘅「修煉」，相信你已經由一個「數字新手」成長為一名相當從容嘅手機用戶喇。但科技生活總有輪迴——新歡終會變舊愛。當你準備等老夥計「退役」時，你知點樣體面、安全咁告別嗎？}

\booksection{一、舊手機嘅「安樂死」：安全出售或轉贈指南}

核心原則：數據清理唔係「刪除」，而係「覆蓋」。好多人以為將相片拖入資源回收筒、清空相冊就萬事大吉，結果手機轉手後，新主人用數據恢復軟件一掃，舊相片、聊天記錄、身份證相片全部返曬嚟——呢種「社死現場」每年都喺度上演。

\begin{enumerate}
  \item \textbf{自我「大掃除」}：完成最後一次完整備份，退出所有帳號（iCloud、華為/小米/三星帳號、Google帳號），關閉「查找我嘅手機」功能——否則下一位用戶會被激活鎖攔喺門外。
  \item \textbf{執行終極「格式化」}：現代手機嘅「恢復出廠設定」通常包含加密擦除——佢銷毀嘅係加密密鑰本身，密鑰一冇，整盤數據瞬間變為亂碼。
  \item \textbf{最後嘅「障眼法」（僅舊設備需要）}：現代手機（iOS 8+、Android 6+）已預設全盤加密（FBE），恢復出廠設定會銷毀加密密鑰，數據理論上唔可恢復，\textbf{無需此步驟}。僅對極老舊嘅非加密設備（2015年前），先至需要喺恢復出廠設定後用垃圾檔案填滿儲存，再執行一次恢復出廠設定。
\end{enumerate}

\textbf{各品牌數據擦除路徑速查表}：雖然原理相通，但每個品牌將「恢復出廠設定」收埋嘅位置略有唔同，關鍵差異如下。

\begin{center}
\begin{tabularx}{\linewidth}{|X|X|X|}
\hline
\textbf{品牌} & \textbf{擦除入口} & \textbf{關鍵差異與注意事項} \tabularnewline
\hline
華為 & 設定 > 系統同更新 > 重置 > 恢復出廠設定 & 需先退出華為帳號並關閉「查找我嘅手機」，否則會觸發激活鎖；純淨模式下可能需要密碼驗證。 \tabularnewline
小米 / Redmi & 設定 > 我嘅設備 > 恢復出廠設定 & 底部有「格式化虛擬SD卡」同「清除全部內容與設定」選項，兩者都要勾選；退出小米帳號前先關閉「查找設備」。 \tabularnewline
OPPO / 一加 & 設定 > 其他設定 > 還原手機 > 徹底清除全部數據 & ColorOS 提供「徹底清除」模式，會多次覆蓋儲存；需輸入鎖屏密碼確認。 \tabularnewline
vivo & 設定 > 系統管理 > 備份與重置 > 清除所有數據 & 勾選「格式化手機儲存」先至會清理SD卡區；需退出vivo帳號，關閉「查找手機」。 \tabularnewline
三星 & 設定 > 通用管理 > 重置 > 恢復出廠設定 & 需先移除Samsung帳號並關閉「查找我嘅手機」；部分機型有「自動覆蓋」選項，建議開啟。 \tabularnewline
iPhone & 設定 > 通用 > 傳輸或還原iPhone > 抹掉所有內容同設定 & 必須先退出Apple ID，否則下一位用戶被iCloud激活鎖攔住；iOS 15+支援「退出時保留iCloud數據」以備後查。 \tabularnewline
\hline
\end{tabularx}
\end{center}

\begin{tcbwarn}{最容易被忽略嘅三件事}
\begin{itemize}
  \item \textbf{未退出帳號}：手機賣畀下家後被遠程鎖定，自己都落唔到，最後淨係可以揸身份證去官方售後解鎖。
  \item \textbf{SD卡未取出}：恢復出廠設定唔會格式化外置SD卡入面便嘅相片，記得手動取出或單獨格式化。
  \item \textbf{USB手指/OTG外設未斷開}：擦除過程中如果連住外設，可能誤清外設數據，亦可能導致擦除中斷。
\end{itemize}

\end{tcbwarn}

\booksection{二、數字遺產：你嘅社交靈魂點樣安放？}

手機入面便嘅數據，早就唔係「虛擬」嘅喇。佢承載住你嘅資金、社交關係、創作成果，乃至你嘅「數字身份」。\textbf{一個人去世後，呢啲資產點樣處理？}呢個唔再係一個抽象問題，而係越來越多家庭真實面臨嘅難題——親屬唔知密碼、平臺唔畀繼承、加密貨幣永遠丟失。

\begin{tcbinfo}{法律依據：個人資訊保護法第49條}
2021年施行嘅《中華人民共同國個人資訊保護法》第49條明確規定：\textbf{自然人身故嘅，其近親屬為咗自身嘅合法、正當利益，可以對死者相關個人資訊行使本章規定嘅查閲、複製、更正、刪除等權利；死者生前另有安排嘅除外。}

呢個意味住：

\begin{itemize}
  \item \textbf{近親屬有法律依據}向平臺申請處理死者嘅帳號同個人資訊。
  \item 「合法、正當利益」通常包括辦理遺產繼承、關閉帳號、刪除敏感私隱等。
  \item 死者生前若有遺囑或協議另行安排（如指定遺產聯絡人），從其約定。
  \item 平臺唔可以以「用戶私隱」為由一概拒絕配合，但具體流程各家唔同。
\end{itemize}

\end{tcbinfo}

\booksubsection{微信同支付寶：帳號唔可以繼承，但資金可以}

呢兩個中國人最常用嘅App，規則各有唔同：

\begin{itemize}
  \item \textbf{微信帳號}：依據《騰訊微信軟件許可及服務協議》，帳號所有權歸騰訊，用戶僅享有使用權，\textbf{帳號本身唔可以繼承、唔可以轉讓}。但帳號內嘅\textbf{零錢、零錢通資金}屬於個人財產，近親屬可憑死亡證明、親屬關係證明、繼承公證書向財付通申請提現或轉移。
  \item \textbf{支付寶}：帳號同樣唔可以直接繼承，但支付寶餘額、餘額寶、理財通等\textbf{資金屬於可繼承財產}。繼承人可透過支付寶「客服熱線」或線下提交死亡證明、繼承公證書辦理；螞蟻財富類理財產品需走基金公司流程。
  \item \textbf{聊天記錄}：法律上屬個人私隱，平臺唔會主動提供畀家屬。如涉及遺產線索，可透過訴訟或公證申請調取。
\end{itemize}

\begin{tcbwarn}{千萬唔好做嘅事}
\begin{itemize}
  \item \textbf{唔好將死者帳號密碼直接交畀家人「接住用」}——呢個違反平臺協議，資金被凍結甚至觸發風控時舉證更難。
  \item \textbf{唔好嘗試用死者手機收驗證碼「提現」}——可能被認定為盜用，反而觸犯法律。
  \item \textbf{唔好拖延}：部分帳號長期未登入會被自動註銷，資金搵返會更麻煩。
\end{itemize}

\end{tcbwarn}

\booksubsection{Apple「遺產聯絡人」：提前指定，事半功倍}

iOS 15.2 起蘋果推出「遺產聯絡人」功能，允許用戶提前指定信任嘅聯絡人，身故後該聯絡人可憑\textbf{死亡證明 + 遺產聯絡人訪問密鑰}申請訪問死者嘅iCloud數據（相片、備忘錄、通訊錄、備忘錄等）。操作步驟：

\begin{enumerate}
  \item 喺iPhone上進入\textbf{[設定] > [你嘅姓名] > [登入與安全] > [遺產聯絡人]}。
  \item 㩒「添加遺產聯絡人」，使用面容ID確認身份。
  \item 選擇聯絡人（由通訊錄挑選），系統會生成一份\textbf{訪問密鑰}。
  \item 將訪問密鑰透過短訊、郵件或列印紙本交畀對方，並提醒對方妥善保存。
  \item 可添加多個聯絡人作為備份，避免單點失效。
\end{enumerate}

被指定人唔使Apple設備，憑密鑰同死亡證明即可喺蘋果官網申請訪問。整個流程通常喺7日內完成。

\booksubsection{Google「不活躍賬戶管理器」：自動處理身後事}

Google提供「不活躍賬戶管理器（Inactive Account Manager）」，能喺用戶長期未登入時自動執行預設操作——呢個係目前最完善嘅數字遺產機制之一。

\begin{enumerate}
  \item 訪問\textbf{myaccount.google.com/inactive}，登入Google帳號。
  \item 設定超時期限：3 / 6 / 12 / 18個月未登入後觸發。
  \item 添加受信任聯絡人（最多10人），並設定每位可訪問邊啲數據（Gmail、相冊、雲端硬碟等）。
  \item 可選：設定自動回覆郵件，告知來信者你已無法處理。
  \item 最後選擇：到期後係\textbf{分享數據畀聯絡人}，仲係\textbf{直接刪除賬戶}。
  \item 設定手機號作為兜底驗證，避免誤觸發。
\end{enumerate}

\booksubsection{國內社交帳號：紀念化與註銷流程}

國內主流社交平臺對逝者帳號嘅處理分兩種思路：\textbf{「紀念化」（保留帳號等親友緬懷）}同\textbf{「註銷」（徹底刪除）}。家屬可憑死亡證明、親屬關係證明、申請人身份證向平臺客服申請。

\begin{itemize}
  \item \textbf{微博}：支援帳號「紀念化」，由親屬提供死亡證明申請，帳號會被鎖定無法登入但內容保留，主頁會顯示「該帳號已設定為紀念帳號」。亦可申請註銷，需提供繼承公證。
  \item \textbf{抖音}：透過「我 - 設定 - 反饋與幫助 - 帳號與安全 - 帳號註銷」或人工客服，提交死亡證明、親屬關係證明、申請人身份證，審核通過後可註銷或轉為紀念狀態。
  \item \textbf{B站}：客服渠道提交逝者證明及親屬證明，可申請帳號註銷或保留為紀念頁；UP主嘅歷史稿件處理需結合版權聲明。
  \item \textbf{小紅書}：透過官方客服申訴，提交死亡證明、親屬關係證明，可申請註銷帳號並刪除內容。
\end{itemize}

\begin{tcbinfo}{關鍵提醒}
\begin{itemize}
  \item 各平臺規則會更新，辦理前先查閲官方客服最新指引。
  \item 所有證明檔案建議先掃描備份再提交，原件妥善保管。
  \item 跨平臺一次性辦理時，可統一申請一份\textbf{繼承公證書}，慳返跑幾次手續。
\end{itemize}

\end{tcbinfo}

\booksubsection{數字資產嘅繼承難題}

除咗社交帳號，越來越多人持有\textbf{真正有價值嘅數字資產}，其繼承問題更加複雜：

\begin{itemize}
  \item \textbf{加密貨幣（比特幣、以太坊等）}：基於「私鑰即所有權」原則，\textbf{冇私鑰就等於資產消失}。繼承人無法透過法律檔案向鏈上「申訴」。建議：將助記詞/私鑰分片存放喺多個物理安全位置（如銀行保險箱+屋企保險櫃），並喺遺囑中註明訪問方式。
  \item \textbf{交易所賬戶}：同銀行賬戶類似，可憑繼承公證書向交易所（如幣安、OKX）申請處置，但跨境交易所流程漫長。
  \item \textbf{域名}：屬於可繼承財產，繼承人喺域名註冊商（如阿里雲、GoDaddy）提交繼承公證後可辦理過戶。
  \item \textbf{知識產權（原創文章、影片、專利、軟件著作權）}：著作權中嘅財產權利可繼承（人身權唔繼承），收益權可由繼承人享有至作者身故後50年。
  \item \textbf{雲端儲存與網盤}：內容數據歸屬個人，可繼承；但部分平臺協議規定帳號本身唔可轉讓，需先下載再註銷。
\end{itemize}

\textbf{加密貨幣「消失」嘅真實代價}\par

據Chainalysis報告，全球已有\textbf{超過300萬枚比特幣}因私鑰丟失而永遠無法搵返，相當於現有總量嘅六分之一。如果你持有加密資產，請務必等至少一位你絕對信任嘅人知道\textbf{恢復路徑}——否則呢啲財富將永遠沉睡喺區塊鏈上。

\booksubsection{一份「數字遺囑」模板}

傳統嘅財產遺囑管唔到數字資產。建議你每年更新一份「數字遺囑」附件，同正式遺囑一併存放。以下模板供參考：

\textbf{《個人數字資產處置附件》（模板）}\par

\textbf{第一部分：帳號清單}（僅記錄帳號名與所屬平臺，\textbf{唔寫密碼}）

\begin{itemize}
  \item 社交類：微信（帳號xxx）、微博、抖音、B站、小紅書等
  \item 支付類：支付寶、微信支付、各銀行卡App
  \item 雲端類：iCloud、華為雲、小米雲、Google帳號、微軟帳號
  \item 郵箱類：QQ郵箱、網易郵箱、Gmail等
  \item 其他：網盤、訂閲服務、遊戲帳號等
\end{itemize}

\textbf{第二部分：密碼託管方案}

\begin{itemize}
  \item 主密碼存放方式：使用密碼管理器（如Bitwarden、1Password），主密碼寫喺紙本中，封存於家中保險櫃。
  \item 恢復密鑰/助記詞：分片存放，至少兩處物理位置（如銀行保管箱+家中保險櫃）。
  \item 授權查詢人：姓名、身份證號、與本人關係、聯繫方式。
\end{itemize}

\textbf{第三部分：處置意願}

\begin{itemize}
  \item 各類帳號希望「紀念化」定係「註銷」。
  \item iCloud / Google 數據係咪指定「遺產聯絡人」或「不活躍賬戶管理器」。
  \item 加密資產由邊個繼承，私鑰獲取路徑說明。
  \item 版權類作品收益由邊個享有。
\end{itemize}

\textbf{第四部分：聯絡人與平臺}

\begin{itemize}
  \item 信任執行人姓名、聯繫方式、備用聯絡人。
  \item 常駐平臺客服入口與所需證明材料清單。
  \item 本附件更新日期：\_\_\_\_年\_\_月\_\_日，下次複核日期：\_\_\_\_年\_\_月\_\_日。
\end{itemize}

\begin{tcbwarn}{關於「密碼」嘅法律提示}
將密碼直接寫入遺囑存在洩露風險，且遺囑公示後所有繼承人都能睇到。\textbf{更穩妥嘅做法係：遺囑入面便淨係寫「帳號清單 + 處置意願 + 密碼託管位置」，真正嘅密碼透過密碼管理器主密碼單獨傳畀信任執行人}。咁樣既符合法律要求，亦降低咗密碼洩露風險。

\end{tcbwarn}

\booksection{三、環保嘅最後一程}

一部手機入面便含金、銀、銅、稀土等數十種金屬，隨意丟棄既污染土壤，亦浪費資源。正確嘅回收渠道有三類，各有側重：

\begin{center}
\begin{tabularx}{\linewidth}{|X|X|X|X|}
\hline
\textbf{渠道} & \textbf{回收價} & \textbf{數據安全} & \textbf{適用場景} \tabularnewline
\hline
官方以舊換新（華為/蘋果/小米等） & 中等，常疊加補貼券 & 最高，官方有規範流程 & 正打算換同品牌新機時最劃算 \tabularnewline
第三方回收平臺（愛回收、轉轉、回收寶等） & 略高，可比價 & 較高，有數據清除報告 & 想直接變現、唔綁定品牌 \tabularnewline
公益捐贈（如鄉村學校、愛心機構） & 無經濟回報 & 需自行徹底擦除 & 設備仍可用、想做善事 \tabularnewline
發揮餘熱（自留改造） & 慳返新設備錢 & 原機主本人繼續用，無風險 & 設備效能尚可、屋企有改造需求 \tabularnewline
\hline
\end{tabularx}
\end{center}

\begin{itemize}
  \item \textbf{官方回收/以舊換新}：最推薦嘅方式，流程規範，數據安全有保障。蘋果、華為、小米等品牌官網都有「以舊換新」入口，評估後可抵扣新機款。
  \item \textbf{正規回收平臺}：選擇信譽良好嘅大型回收平臺，落單後順豐上門取件，檢測後打款。務必選擇\textbf{提供「數據清除報告」嘅}平臺，並保留憑證。
  \item \textbf{公益捐贈}：可聯絡當地鄉村學校、社區公益組織，舊手機可改裝成教學平板或老人機使用。\textbf{捐贈前務必執行完整嘅擦除流程，並卸除所有SIM/SD卡}。
  \item \textbf{發揮餘熱}：改成車載導航儀、智能家居中控、監控鏡頭、BB監視器、電子相框。一部舊手機可以再戰好幾年。
\end{itemize}

\begin{tcbinfo}{電子垃圾嘅真相}
一部手機嘅電路板入面便含有\textbf{金、銀、鈀、銅}等貴金屬，每噸廢舊電路板含金量係金礦石嘅數十倍。但若隨意丟棄，其中嘅鉛、汞、鎘會滲入土壤與地下水，污染持續數十年。\textbf{正確回收，既保護環境，亦能等稀缺金屬循環再生。}

\end{tcbinfo}

\begin{tcbpractice}{實操：為你嘅舊手機執行一場「標準葬禮」}
\begin{enumerate}
  \item 最後一次完整備份：相片上傳雲端、重要檔案複製到電腦、聊天記錄匯出存檔。
  \item 退出所有帳號：iCloud / 華為 / 小米 / Google 帳號，關閉「查找我嘅手機」。
  \item 取出SIM卡與SD卡，斷開所有OTG外設。
  \item 執行加密恢復出廠設定，必要時手動勾選「格式化儲存」。
  \item （僅舊設備）用垃圾檔案填滿儲存，再恢復出廠設定一次。
  \item 選擇佢嘅下一站：官方以舊換新、正規回收平臺、公益捐贈，或轉贈家人。
  \item 同步更新你嘅「數字遺產清單」，將已淘汰嘅設備由清單中劃走。
\end{enumerate}

\end{tcbpractice}

\begin{tcbpractice}{本月任務清單}
\begin{enumerate}
  \item 為你目前嘅主力手機指定Apple「遺產聯絡人」或設定Google「不活躍賬戶管理器」。
  \item 用本節模板起草一份《個人數字資產處置附件》，同家人溝通存放位置。
  \item 檢查屋企係咪有2部以上閒置舊手機，預約一次官方或正規平臺回收。
\end{enumerate}

\end{tcbpractice}
